\documentclass[11pt, a4paper]{article}

%% >>> preamble
\usepackage[a4paper, margin=2.2cm]{geometry}
\usepackage{amsmath, amssymb, amsthm}
\usepackage{enumitem}
\usepackage{fancyhdr}
\usepackage{tcolorbox}
\pagestyle{fancy}
\fancyhf{}
\lhead{\textbf{MATH 101} -- Problem Set 1}
\rhead{Your Name}
\cfoot{\thepage}
\setlength{\headheight}{14pt}
\newcounter{problem}
\newenvironment{problem}[1][]{%
  \refstepcounter{problem}%
  \begin{tcolorbox}[colback=gray!5, colframe=gray!55, boxrule=0.5pt, arc=1mm,
      colbacktitle=gray!18, coltitle=black, fonttitle=\bfseries,
      title={Problem \theproblem\quad #1}]%
}{\end{tcolorbox}}
\newenvironment{solution}{\begin{proof}[Solution]}{\end{proof}}
%% <<< preamble

\begin{document}
%% >>> snippet
\begin{center}
  {\LARGE\bfseries Problem Set 1}\\[4pt]
  Due Friday, 11:59\,pm
\end{center}

\begin{problem}[10 points]
  Show that $\displaystyle\sum_{k=1}^{n} k = \frac{n(n+1)}{2}$ for all $n \in \mathbb{N}$.
\end{problem}
\begin{solution}
  By induction on $n$. For $n = 1$ both sides equal $1$. Assuming the claim for $n$,
  \[
    \sum_{k=1}^{n+1} k = \frac{n(n+1)}{2} + (n + 1) = \frac{(n+1)(n+2)}{2},
  \]
  which is the claim for $n + 1$.
\end{solution}

\begin{problem}[5 points]
  Solve for $x$:
  \begin{enumerate}[label=(\alph*)]
    \item $2x + 3 = 11$,
    \item $x^2 - 5x + 6 = 0$.
  \end{enumerate}
\end{problem}
\begin{solution}
  \begin{enumerate}[label=(\alph*)]
    \item $x = 4$.
    \item $(x - 2)(x - 3) = 0$, so $x \in \{2, 3\}$.
  \end{enumerate}
\end{solution}
%% <<< snippet
\end{document}
